% GENERATED FILE. Edit canonical lessons, then run npm run generate:books.
% canonical-source-hash: fnv1a64:04d6216d58a04d5b
% canonical-lessons: HI-C465-mazaq, HI-C465-hasana, HI-C465-rulana, HI-C465-ah, HI-C465-vilap

\chapter{Laughter and Tears}
\label{ch:hi-laughter-and-tears}
\hlchaptermodality{465}

\begin{chapteropening}
\textbf{By the end of this chapter:} \emph{I can say what made someone laugh or cry.}

Complete the last lesson of chapter 465: I can say what made someone laugh or cry.
\end{chapteropening}

\section[\texorpdfstring{\hi{मज़ाक़} (mazāq)}{मज़ाक़ (mazāq)}]{\hi{मज़ाक़} (mazāq) — a joke; teasing}

\label{lesson:HI-C465-mazaq}



\begin{quote}
Before the new one: say the Hindi for curiosity, then the Hindi for interesting.
\end{quote}

\subsection*{You'll want to know: \hi{मज़ाक़}}
\textbf{\hi{मज़ाक़}} — \emph{mazāq} — \textquotedblleft{}a joke; teasing\textquotedblright{}.

\begin{quote}
\hi{अरुण} \hi{के} \hi{मज़ाक़} \hi{पर} \hi{सब} \hi{हँसे।}
\end{quote}

\begin{quote}
\emph{Aruṇ ke mazāq par sab hãse.}
\end{quote}

\begin{quote}
Everyone laughed at Arun's joke.
\end{quote}

Masculine.

\begin{cousinweb}
\textbf{\hi{मज़ाक़}} (\emph{mazāq}) is usually traced to Arabic \emph{madhāq}, \textquotedblleft{}taste, flavour\textquotedblright{}, which reached Hindi through Persian, where it came to mean a taste for fun and then a joke.
\end{cousinweb}

\begin{grammarlens}[title={What you've built}]
One more word for describing what you lived and felt.
\end{grammarlens}

\subsection*{Your turn}
\begin{itemize}
\raggedright
  \item \emph{Say it:} \emph{mazāq}
  \item \emph{Say it:} \emph{mazāq}, once more
  \item \emph{Say it:} \emph{Aruṇ ke mazāq par sab hãse.} — Everyone laughed at Arun's joke.
  \item \emph{Recall:} say the Hindi for curiosity, then the Hindi for interesting, then say \emph{mazāq} again
\end{itemize}

\begin{tcolorbox}[breakable,colback=teal!4,colframe=teal!35!black,title={Before you move on}]
What does \hi{मज़ाक़} mean? (\textquotedblleft{}A joke; teasing\textquotedblright{}.) Say it once more.
\end{tcolorbox}
\section[\texorpdfstring{\hi{हँसाना} (hãsānā)}{हँसाना (hãsānā)}]{\hi{हँसाना} (hãsānā) — to make someone laugh}

\label{lesson:HI-C465-hasana}



\begin{quote}
Before the new one: say the Hindi for interesting, then the Hindi for a joke.
\end{quote}

\subsection*{You'll want to know: \hi{हँसाना}}
\textbf{\hi{हँसाना}} — \emph{hãsānā} — \textquotedblleft{}to make someone laugh\textquotedblright{}.

\begin{quote}
\hi{अरुण} \hi{अपने} \hi{मज़ाक़} \hi{से} \hi{सब} \hi{लोगों} \hi{को} \hi{हँसाता} \hi{है।}
\end{quote}

\begin{quote}
\emph{Aruṇ apne mazāq se sab logoṁ ko hãsātā hai.}
\end{quote}

\begin{quote}
Arun makes everyone laugh with his jokes.
\end{quote}

Past, with \hi{ने} (he or she …): \hi{उसने} \hi{बच्चों} \hi{को} \hi{हँसाया}, \emph{usne baccoṁ ko hãsāyā}.

\begin{cousinweb}
\textbf{\hi{हँसाना}} (\emph{hãsānā}) is the causative partner of \emph{hãsnā}, \textquotedblleft{}to laugh\textquotedblright{}, which CDIAL lists under Sanskrit \emph{hásati}, \textquotedblleft{}laughs\textquotedblright{} (CDIAL 14021).
\end{cousinweb}

\begin{grammarlens}[title={What you've built}]
One more word for describing what you lived and felt.
\end{grammarlens}

\subsection*{Your turn}
\begin{itemize}
\raggedright
  \item \emph{Say it:} \emph{hãsānā}
  \item \emph{Say it:} \emph{hãsānā}, once more
  \item \emph{Say it:} \emph{Aruṇ apne mazāq se sab logoṁ ko hãsātā hai.} — Arun makes everyone laugh with his jokes.
  \item \emph{Recall:} say the Hindi for interesting, then the Hindi for a joke, then say \emph{hãsānā} again
\end{itemize}

\begin{tcolorbox}[breakable,colback=teal!4,colframe=teal!35!black,title={Before you move on}]
What does \hi{हँसाना} mean? (\textquotedblleft{}To make someone laugh\textquotedblright{}.) Say it once more.
\end{tcolorbox}
\section[\texorpdfstring{\hi{रुलाना} (rulānā)}{रुलाना (rulānā)}]{\hi{रुलाना} (rulānā) — to make someone cry}

\label{lesson:HI-C465-rulana}



\begin{quote}
Before the new one: say the Hindi for a joke, then the Hindi for to make someone laugh.
\end{quote}

\subsection*{You'll want to know: \hi{रुलाना}}
\textbf{\hi{रुलाना}} — \emph{rulānā} — \textquotedblleft{}to make someone cry\textquotedblright{}.

\begin{quote}
\hi{दादी} \hi{की} \hi{पुरानी} \hi{कहानी} \hi{ने} \hi{मीरा} \hi{को} \hi{रुला} \hi{दिया।}
\end{quote}

\begin{quote}
\emph{dādī kī purānī kahānī ne Mīrā ko rulā diyā.}
\end{quote}

\begin{quote}
Grandmother's old story moved Mira to tears.
\end{quote}

Past, with \hi{ने} (he or she …): \hi{उसने} \hi{बच्चे} \hi{को} \hi{रुलाया}, \emph{usne bacce ko rulāyā}.

\begin{cousinweb}
\textbf{\hi{रुलाना}} (\emph{rulānā}) is the causative partner of \emph{ronā}, \textquotedblleft{}to weep\textquotedblright{}, which CDIAL lists under Sanskrit \emph{rōdati}, \textquotedblleft{}weeps\textquotedblright{} (CDIAL 10840); the vowel shortens before the causative ending, as it does in \emph{sulānā} beside \emph{sonā}.
\end{cousinweb}

\begin{grammarlens}[title={What you've built}]
One more word for describing what you lived and felt.
\end{grammarlens}

\subsection*{Your turn}
\begin{itemize}
\raggedright
  \item \emph{Say it:} \emph{rulānā}
  \item \emph{Say it:} \emph{rulānā}, once more
  \item \emph{Say it:} \emph{dādī kī purānī kahānī ne Mīrā ko rulā diyā.} — Grandmother's old story moved Mira to tears.
  \item \emph{Recall:} say the Hindi for a joke, then the Hindi for to make someone laugh, then say \emph{rulānā} again
\end{itemize}

\begin{tcolorbox}[breakable,colback=teal!4,colframe=teal!35!black,title={Before you move on}]
What does \hi{रुलाना} mean? (\textquotedblleft{}To make someone cry\textquotedblright{}.) Say it once more.
\end{tcolorbox}
\section[\texorpdfstring{\hi{आह} (āh)}{आह (āh)}]{\hi{आह} (āh) — a sigh}

\label{lesson:HI-C465-ah}



\begin{quote}
Before the new one: say the Hindi for to make someone laugh, then the Hindi for to make someone cry.
\end{quote}

\subsection*{You'll want to know: \hi{आह}}
\textbf{\hi{आह}} — \emph{āh} — \textquotedblleft{}a sigh\textquotedblright{}.

\begin{quote}
\hi{थककर} \hi{दादी} \hi{ने} \hi{एक} \hi{लंबी} \hi{आह} \hi{भरी।}
\end{quote}

\begin{quote}
\emph{thakkar dādī ne ek lambī āh bharī.}
\end{quote}

\begin{quote}
Worn out, my grandmother heaved a long sigh.
\end{quote}

Feminine.

\begin{cousinweb}
\textbf{\hi{आह}} (\emph{āh}) is the sound of a sigh made into a noun; Monier-Williams gives Sanskrit \emph{āha} as an exclamation (Monier-Williams). \emph{āh bharnā}, \textquotedblleft{}to fill a sigh\textquotedblright{}, is to heave one.
\end{cousinweb}

\begin{grammarlens}[title={What you've built}]
One more word for describing what you lived and felt.
\end{grammarlens}

\subsection*{Your turn}
\begin{itemize}
\raggedright
  \item \emph{Say it:} \emph{āh}
  \item \emph{Say it:} \emph{āh}, once more
  \item \emph{Say it:} \emph{thakkar dādī ne ek lambī āh bharī.} — Worn out, my grandmother heaved a long sigh.
  \item \emph{Recall:} say the Hindi for to make someone laugh, then the Hindi for to make someone cry, then say \emph{āh} again
\end{itemize}

\begin{tcolorbox}[breakable,colback=teal!4,colframe=teal!35!black,title={Before you move on}]
What does \hi{आह} mean? (\textquotedblleft{}A sigh\textquotedblright{}.) Say it once more.
\end{tcolorbox}
\section[\texorpdfstring{\hi{विलाप} (vilāp)}{विलाप (vilāp)}]{\hi{विलाप} (vilāp) — wailing; lamentation}

\label{lesson:HI-C465-vilap}



\begin{quote}
Before the new one: say the Hindi for to make someone cry, then the Hindi for a sigh.
\end{quote}

\subsection*{You'll want to know: \hi{विलाप}}
\textbf{\hi{विलाप}} — \emph{vilāp} — \textquotedblleft{}wailing; lamentation\textquotedblright{}.

\begin{quote}
\hi{बेटी} \hi{का} \hi{विलाप} \hi{सुनकर} \hi{सब} \hi{लोग} \hi{रो} \hi{पड़े।}
\end{quote}

\begin{quote}
\emph{beṭī kā vilāp sunkar sab log ro paṛe.}
\end{quote}

\begin{quote}
Hearing the daughter's wailing, everyone broke down in tears.
\end{quote}

Masculine.

\begin{cousinweb}
\textbf{\hi{विलाप}} (\emph{vilāp}) is a tatsama, borrowed from Sanskrit \emph{vilāpa}, \textquotedblleft{}lamentation\textquotedblright{}, from \emph{vi-lap}, \textquotedblleft{}to wail\textquotedblright{} (Monier-Williams).
\end{cousinweb}

\begin{grammarlens}[title={What you've built}]
One more word for describing what you lived and felt.
\end{grammarlens}

\subsection*{Your turn}
\begin{itemize}
\raggedright
  \item \emph{Say it:} \emph{vilāp}
  \item \emph{Say it:} \emph{vilāp}, once more
  \item \emph{Say it:} \emph{beṭī kā vilāp sunkar sab log ro paṛe.} — Hearing the daughter's wailing, everyone broke down in tears.
  \item \emph{Recall:} say the Hindi for to make someone cry, then the Hindi for a sigh, then say \emph{vilāp} again
\end{itemize}

\begin{tcolorbox}[breakable,colback=teal!4,colframe=teal!35!black,title={Before you move on}]
What does \hi{विलाप} mean? (\textquotedblleft{}Wailing; lamentation\textquotedblright{}.) Say it once more.
\end{tcolorbox}
